% figures/w01-otter-layout.svg — Otter 를 위에서 본 그림, 그리고 B 가 왜 그 모양인가.
%
%   bash _tools/tikz2svg.sh figures/src/w01-otter-layout.tex figures/w01-otter-layout.svg
%
% 치수는 otter.m / _tools/otter_config.m 의 것이다:
%   y_pont = 0.395 m,  B_pont = 0.25 m,  L = 2.0 m.
% 갑판 평면도이므로 x_b(선수)를 위로, y_b(우현)를 오른쪽으로 둔다.
%
% B 의 값은 눈대중이 아니라 _tools/otter_B.m 의 열 규칙에서 나온다:
%
%     열 = [ e_x ; e_y ; x e_y - y e_x ]
%
% 'base' 형상은 pos = [0 0; -0.395 +0.395], dir = [1 1; 0 0] 이므로
%     좌현 열 = [1; 0; 0 - (-0.395)(1)] = [1; 0; +0.395]
%     우현 열 = [1; 0; 0 - (+0.395)(1)] = [1; 0; -0.395]
% 세 번째 성분의 부호가 y 좌표와 **반대**라는 것이 이 그림의 요점 하나다.
\documentclass[border=5pt]{standalone}
% gnc-style.tex — 이 볼트의 TikZ 그림이 공유하는 스타일.
%
% 저널 그림 규약을 스타일 하나로 굳혀 둔다. 그림마다 선 굵기나 화살촉을
% 다시 정하지 않는다 — 그것이 그림이 제각각으로 보이는 원인이다.
%
% 쓰는 법:  % gnc-style.tex — 이 볼트의 TikZ 그림이 공유하는 스타일.
%
% 저널 그림 규약을 스타일 하나로 굳혀 둔다. 그림마다 선 굵기나 화살촉을
% 다시 정하지 않는다 — 그것이 그림이 제각각으로 보이는 원인이다.
%
% 쓰는 법:  % gnc-style.tex — 이 볼트의 TikZ 그림이 공유하는 스타일.
%
% 저널 그림 규약을 스타일 하나로 굳혀 둔다. 그림마다 선 굵기나 화살촉을
% 다시 정하지 않는다 — 그것이 그림이 제각각으로 보이는 원인이다.
%
% 쓰는 법:  \input{gnc-style}  (figures/src/ 안에서)

\usepackage{tikz}
\usepackage{amsmath}
\usepackage{xcolor}
\usetikzlibrary{arrows.meta, positioning, calc, fit, backgrounds, decorations.pathmorphing}

% ---- 색 --------------------------------------------------------------------
% 색은 강조 하나에만 쓴다. 나머지는 검정.
\definecolor{ink}{HTML}{1A1A1A}
\definecolor{accent}{HTML}{7C3AED}
\definecolor{warn}{HTML}{B91C1C}
\definecolor{soft}{HTML}{666666}

% ---- 선과 화살촉 -----------------------------------------------------------
\tikzset{
  % 신호선. 화살촉은 작고 뾰족한 것 하나뿐이다.
  sig/.style   = {ink, line width=0.5pt, -{Stealth[length=4pt, width=3pt]}},
  wire/.style  = {ink, line width=0.5pt},
  % 강조 경로 — 파선으로 구분한다. 흑백 인쇄에서도 살아야 하기 때문이다.
  asig/.style  = {accent, line width=0.5pt, densely dashed,
                  -{Stealth[length=4pt, width=3pt]}},
  awire/.style = {accent, line width=0.5pt, densely dashed},
  % 블록. 흰 바탕, 직각, 검은 테두리.
  blk/.style   = {draw=ink, line width=0.55pt, fill=white,
                  minimum width=17mm, minimum height=8mm, inner sep=2pt, font=\small},
  % 합산점
  sum/.style   = {draw=ink, line width=0.55pt, fill=white, circle,
                  inner sep=0pt, minimum size=4.6mm, font=\scriptsize},
  asum/.style  = {draw=accent, line width=0.55pt, fill=white, circle,
                  inner sep=0pt, minimum size=4.6mm, font=\scriptsize},
  % 분기점
  dot/.style   = {ink, fill=ink, circle, inner sep=0pt, minimum size=1.5mm},
  % 신호 이름 — 선 위에 작게
  sname/.style = {font=\small, inner sep=1.5pt},
  % 그림 안의 짧은 설명. 그림 안의 글은 최소로 둔다
  note/.style  = {font=\scriptsize, soft, inner sep=1.5pt},
  anote/.style = {font=\scriptsize, accent, inner sep=1.5pt},
  % 축
  axis/.style  = {ink, line width=0.5pt},
  tick/.style  = {font=\scriptsize, soft},
}

% 캡션 한 줄 — 그림 아래 가는 선과 함께
\newcommand{\gnccaption}[2]{%
  \node[anchor=north west, text width=#1, align=left, font=\scriptsize, ink]
        at (current bounding box.south west) {\vspace{2pt}#2};}
  (figures/src/ 안에서)

\usepackage{tikz}
\usepackage{amsmath}
\usepackage{xcolor}
\usetikzlibrary{arrows.meta, positioning, calc, fit, backgrounds, decorations.pathmorphing}

% ---- 색 --------------------------------------------------------------------
% 색은 강조 하나에만 쓴다. 나머지는 검정.
\definecolor{ink}{HTML}{1A1A1A}
\definecolor{accent}{HTML}{7C3AED}
\definecolor{warn}{HTML}{B91C1C}
\definecolor{soft}{HTML}{666666}

% ---- 선과 화살촉 -----------------------------------------------------------
\tikzset{
  % 신호선. 화살촉은 작고 뾰족한 것 하나뿐이다.
  sig/.style   = {ink, line width=0.5pt, -{Stealth[length=4pt, width=3pt]}},
  wire/.style  = {ink, line width=0.5pt},
  % 강조 경로 — 파선으로 구분한다. 흑백 인쇄에서도 살아야 하기 때문이다.
  asig/.style  = {accent, line width=0.5pt, densely dashed,
                  -{Stealth[length=4pt, width=3pt]}},
  awire/.style = {accent, line width=0.5pt, densely dashed},
  % 블록. 흰 바탕, 직각, 검은 테두리.
  blk/.style   = {draw=ink, line width=0.55pt, fill=white,
                  minimum width=17mm, minimum height=8mm, inner sep=2pt, font=\small},
  % 합산점
  sum/.style   = {draw=ink, line width=0.55pt, fill=white, circle,
                  inner sep=0pt, minimum size=4.6mm, font=\scriptsize},
  asum/.style  = {draw=accent, line width=0.55pt, fill=white, circle,
                  inner sep=0pt, minimum size=4.6mm, font=\scriptsize},
  % 분기점
  dot/.style   = {ink, fill=ink, circle, inner sep=0pt, minimum size=1.5mm},
  % 신호 이름 — 선 위에 작게
  sname/.style = {font=\small, inner sep=1.5pt},
  % 그림 안의 짧은 설명. 그림 안의 글은 최소로 둔다
  note/.style  = {font=\scriptsize, soft, inner sep=1.5pt},
  anote/.style = {font=\scriptsize, accent, inner sep=1.5pt},
  % 축
  axis/.style  = {ink, line width=0.5pt},
  tick/.style  = {font=\scriptsize, soft},
}

% 캡션 한 줄 — 그림 아래 가는 선과 함께
\newcommand{\gnccaption}[2]{%
  \node[anchor=north west, text width=#1, align=left, font=\scriptsize, ink]
        at (current bounding box.south west) {\vspace{2pt}#2};}
  (figures/src/ 안에서)

\usepackage{tikz}
\usepackage{amsmath}
\usepackage{xcolor}
\usetikzlibrary{arrows.meta, positioning, calc, fit, backgrounds, decorations.pathmorphing}

% ---- 색 --------------------------------------------------------------------
% 색은 강조 하나에만 쓴다. 나머지는 검정.
\definecolor{ink}{HTML}{1A1A1A}
\definecolor{accent}{HTML}{7C3AED}
\definecolor{warn}{HTML}{B91C1C}
\definecolor{soft}{HTML}{666666}

% ---- 선과 화살촉 -----------------------------------------------------------
\tikzset{
  % 신호선. 화살촉은 작고 뾰족한 것 하나뿐이다.
  sig/.style   = {ink, line width=0.5pt, -{Stealth[length=4pt, width=3pt]}},
  wire/.style  = {ink, line width=0.5pt},
  % 강조 경로 — 파선으로 구분한다. 흑백 인쇄에서도 살아야 하기 때문이다.
  asig/.style  = {accent, line width=0.5pt, densely dashed,
                  -{Stealth[length=4pt, width=3pt]}},
  awire/.style = {accent, line width=0.5pt, densely dashed},
  % 블록. 흰 바탕, 직각, 검은 테두리.
  blk/.style   = {draw=ink, line width=0.55pt, fill=white,
                  minimum width=17mm, minimum height=8mm, inner sep=2pt, font=\small},
  % 합산점
  sum/.style   = {draw=ink, line width=0.55pt, fill=white, circle,
                  inner sep=0pt, minimum size=4.6mm, font=\scriptsize},
  asum/.style  = {draw=accent, line width=0.55pt, fill=white, circle,
                  inner sep=0pt, minimum size=4.6mm, font=\scriptsize},
  % 분기점
  dot/.style   = {ink, fill=ink, circle, inner sep=0pt, minimum size=1.5mm},
  % 신호 이름 — 선 위에 작게
  sname/.style = {font=\small, inner sep=1.5pt},
  % 그림 안의 짧은 설명. 그림 안의 글은 최소로 둔다
  note/.style  = {font=\scriptsize, soft, inner sep=1.5pt},
  anote/.style = {font=\scriptsize, accent, inner sep=1.5pt},
  % 축
  axis/.style  = {ink, line width=0.5pt},
  tick/.style  = {font=\scriptsize, soft},
}

% 캡션 한 줄 — 그림 아래 가는 선과 함께
\newcommand{\gnccaption}[2]{%
  \node[anchor=north west, text width=#1, align=left, font=\scriptsize, ink]
        at (current bounding box.south west) {\vspace{2pt}#2};}


\begin{document}
\begin{tikzpicture}[x=1mm, y=1mm]

\def\YP{22.9}        % y_pont = 0.395 m   (1 m = 58 mm)
\def\HB{58}          % 반선장 = 1.0 m
\def\PW{7.25}        % 폰툰 반폭 = 0.125 m

\node[ink, font=\small\bfseries, anchor=north west] at (-40,94)
     {The Otter from above — where the propellers are, and why $\mathbf{B}$ looks the way it does};
\node[anchor=north west, align=left, text width=118mm, font=\scriptsize, soft] at (-40,88) {%
  A catamaran: two pontoons, one propeller in each, both bolted facing forward.
  Every entry of $\mathbf{B}$ follows from this picture and nothing else.};

% ---- 폰툰 둘 ---------------------------------------------------------------
\fill[ink, opacity=0.04] ({-\YP+\PW},-\HB) rectangle ({\YP-\PW},\HB);
\foreach \s in {-1,1} {
  \fill[ink, opacity=0.10] ({\s*\YP-\PW},-\HB) rectangle ({\s*\YP+\PW},\HB);
  \draw[ink, line width=0.5pt, rounded corners=3pt]
        ({\s*\YP-\PW},-\HB) rectangle ({\s*\YP+\PW},\HB);
}

% ---- 중심선 ----------------------------------------------------------------
\draw[soft, line width=0.4pt, densely dashed] (0,{-\HB-8}) -- (0,{\HB+8});
\node[note, anchor=south] at (0,{\HB+9}) {centreline};

% ---- 추진기 둘 -------------------------------------------------------------
% x 위치는 선미로 그렸다. 열 규칙에서 x 는 e_y 와만 곱해지고 e_y = 0 이므로
% 앞을 향한 프로펠러에서는 x 가 B 에 전혀 들어가지 않는다 — 오른쪽 글에 적었다.
\foreach \s in {-1,1} {
  \fill[ink] ({\s*\YP-3.4},-52) rectangle ({\s*\YP+3.4},-42);
  \draw[alng, line width=1.2pt, -{Stealth[length=5pt,width=3.6pt]}]
        ({\s*\YP},-40) -- ({\s*\YP},-14);
}
\node[note, anchor=north] at ({-\YP},{-\HB-2}) {port};
\node[note, anchor=north] at ({ \YP},{-\HB-2}) {starboard};
\node[alng, font=\scriptsize, anchor=south east] at ({-\YP-1},-24) {$T_1$};
\node[alng, font=\scriptsize, anchor=south west] at ({ \YP+1},-24) {$T_2$};

% ---- 몸체 축 ---------------------------------------------------------------
\draw[axis, -{Stealth[length=4.5pt,width=3.4pt]}] (0,0) -- (0,42);
\draw[axis, -{Stealth[length=4.5pt,width=3.4pt]}] (0,0) -- (44,0);
\fill[ink] (0,0) circle (0.8);
\node[font=\small, anchor=east]  at (-1.5,40) {$x_b$};
\node[font=\small, anchor=north] at (43,-1.5) {$y_b$};
\node[font=\scriptsize, anchor=north east] at (-1.5,-1.5) {$o_b$};

% ---- 모멘트 팔 -------------------------------------------------------------
% 브래킷의 끝이 추력선 위에 놓이도록 y = -30 에 그린다.
\draw[accent, line width=0.6pt, |-|] (0,-30) -- ({-\YP},-30);
\draw[accent, line width=0.6pt, |-|] (0,-30) -- ({ \YP},-30);
\node[accent, font=\scriptsize, anchor=north, align=center] at ({-0.5*\YP},-31.5)
     {$y_1$\\ $-0.395$\,m};
\node[accent, font=\scriptsize, anchor=north, align=center] at ({ 0.5*\YP},-31.5)
     {$y_2$\\ $+0.395$\,m};

% ===================== 오른쪽 : B 를 그림에서 읽는다 ========================
\begin{scope}[shift={(92,0)}]
  \node[ink, font=\small\bfseries, anchor=north west] at (0,76)
       {reading $\mathbf{B}$ off the picture};

  \node[anchor=north west, align=left, text width=92mm, font=\scriptsize] at (0,69) {%
    One propeller at $(x_i, y_i)$ pushing along the unit vector
    $\mathbf{e}_i = (e_{x}, e_{y})$ contributes exactly one column
    (\texttt{\_tools/otter\_B.m}, the column rule):};
  \node[anchor=north west, font=\small] at (6,55) {%
    $\mathbf{b}_i = \begin{bmatrix} e_{x}\\ e_{y}\\ x_i e_{y} - y_i e_{x}\end{bmatrix}$};

  \node[anchor=north west, align=left, text width=92mm, font=\scriptsize] at (0,32) {%
    Both propellers are bolted facing forward, so $\mathbf{e}_i = (1,0)$ for both
    and the picture supplies only $y_1 = -0.395$ and $y_2 = +0.395$\,m.\\[4pt]
    \textbf{Surge.} Both push forward: $e_{x} = 1$ twice, so the first row is
    $[\,1\ \ 1\,]$.\\[4pt]
    \textbf{Sway.} Neither has any component across the hull, $e_{y} = 0$, so the
    second row is empty: $Y = 0$ — \textbf{exactly, always}. No shaft speed can
    produce sway.\\[4pt]
    \textbf{Yaw.} The third entry is $-y_i$, which is the \textbf{negative} of the
    coordinate marked on the drawing: the port propeller sits at $y = -0.395$ and
    turns the bow to \emph{starboard}, which is positive yaw.\\[4pt]
    The along-ship station $x_i$ never appears, because it is multiplied by
    $e_{y} = 0$. The propellers may be drawn anywhere along their pontoons and
    $\mathbf{B}$ is unchanged.};

  \node[anchor=north west, font=\small] at (6,-32) {%
    $\mathbf{B} = \begin{bmatrix} 1 & 1\\ 0 & 0\\ 0.395 & -0.395 \end{bmatrix},
     \qquad \boldsymbol\tau = \mathbf{B}\,\mathbf{f}$};

  \node[anchor=north west, align=left, text width=92mm, font=\scriptsize] at (0,-50) {%
    The hull is drawn here at $\psi = 0$, so $\{b\}$ and $\{n\}$ happen to line up.
    At any other heading the two triads differ by $\mathbf{R}(\psi)$ and
    \textbf{$\mathbf{B}$ does not change at all} — it is written entirely in
    $\{b\}$. That is the whole reason the allocation can be designed once and used
    at every heading.};
\end{scope}

\end{tikzpicture}
\end{document}
